\section*{Glossary}

% A
\begin{glossary}
    \textbf{Absolute Value (Moudlus) for \(\mathbb{F}\):} The non-negative value of a real or complex number without regard to its sign.
     For a real number \(a\), it is denoted as \(|a|\), and for a complex number \(z = a + bi\), it is \(\sqrt{a^2 + b^2}\).
\end{glossary}

\begin{glossary}
    \textbf{Absolute Value (Moudlus) for \(\mathbb{M}_n\left(\mathbb{F}\right)\):} For a matrix \(A\), the absolute value \(|A|= (A^*A)^{1/2}\), 
    where \(A^*\) is the conjugate transpose of \(A\).
\end{glossary}

% C
\begin{glossary}
    \textbf{Condition Number:} A measure of the sensitivity of a matrix’s inverse to changes or errors in the input. 
    It is defined as \(cond(A) = ||A|| \cdot ||A^{-1}|| \) or the ratio of the largest to smallest singular values.
\end{glossary}

% D
\begin{glossary}
    \textbf{Determinant:} A scalar-valued function for entries of a square matrix denoted \(det(A) \) in this paper.
\end{glossary}

% E
\begin{glossary}
    \textbf{Eigenvalue:} If \(A\) is a square matrix, then \(\lambda\) is called an eigenvalue if \(A \textbf{x} = A \lambda\) for \(\textbf{x} \in \mathbb{F}^n \setminus \{0\}\).

\end{glossary}

\begin{glossary}
    \textbf{Eigenvector:} If \(A\textbf{x} = A \lambda\) the corresponding \(\textbf{x}\) is the eigenvector of A corresponding to the eigenvalue \(\lambda\).
\end{glossary}

% F
\begin{glossary}
    \textbf{Field:} A set \(F\) which two binary operations, addition and multiplication, which satisfies the field axioms.
\end{glossary}

% H
\begin{glossary}
    \textbf{Hermitian Matrix:} A matrix which is equal to its conjugate transpose, ie \(A = A^*\).
\end{glossary}

% N
\begin{glossary}
    \textbf{Normal Matrix:} A matrix which commutes with its conjugate transpose, ie \(A^*A = AA^*\).
\end{glossary}

% P 
\begin{glossary}
    \textbf{Positive semi-definite Matrix:} Matrix \(A\) is positive semidefinite if any non-zero vector \(\textbf{x}\) 
    the product 
\(\textbf{x}^*A\textbf{x} \geq 0 \). Such matrices have non-negative eigenvalue.
\end{glossary}

% S 
\begin{glossary}
    \textbf{Schur Decomposition:} Takes an arbitrary square matrix \(A\) then it is decomposed into a unitarily simmilar matrix \(U\) and 
    an upper triangular matrix \(T\), so \(A = UTU^*\). 
\end{glossary}

\begin{glossary}
    \textbf{Spectral Decomposition:} Takes an arbitrary square matrix \(A\) and using its eigenvectors to produces a unitary matrix \(P\)
    and eigenvalues to produces diagonal matrix \(D\) so \(A = PDP^*\)
\end{glossary}

% T 
\begin{glossary}
    \textbf{Trace:} The sum of the diagonal elements in a square matrix \(A\).
\end{glossary}